\documentclass[10pt,a4paper]{amsart}
\usepackage[margin=19mm]{geometry}
\usepackage{amsmath,amssymb,mathtools}
\usepackage[scr=rsfs,cal=euler]{mathalfa}
\usepackage{xfrac}
\usepackage[unicode,hidelinks,bookmarks=false]{hyperref}
\newcommand{\ie}{i.e.\ }
\newcommand{\cf}{cf.\ }
\newcommand{\resp}{respectively\ }
\newcommand{\Subsection}{\subsection}
\providecommand{\nobreakdash}{\nobreak\hbox{-}}
\setlength{\emergencystretch}{2em}
\multlinegap0pt
\usepackage{amssymb}
\usepackage{enumerate}
\newtheorem{lemma}{Lemma}
\newtheorem{theorem}{Theorem}
\newcommand{\R}{\mathbb R}
\title{Function spaces and potential theory in the Orlicz setting}
\author{Pablo Ochoa, Ariel Salort}
\date{Source: arXiv:2604.18408v1 (CC BY 4.0)}
\begin{document}
\maketitle

\noindent\emph{Source and licence.} Function spaces and potential theory in the Orlicz setting. Version arXiv:2604.18408v1 (CC BY 4.0), distributed by arXiv under the Creative Commons Attribution 4.0 International licence, http://creativecommons.org/licenses/by/4.0/, as recorded in the arXiv licence metadata for this exact version and shown on the abstract page https://arxiv.org/abs/2604.18408v1. Full licence notice: \texttt{licenses/arxiv-2604.18408-CC-BY-4.0.txt}.\par\smallskip

\noindent\textit{Editorial scope.} Counted unit: the article's theorem teo.s1 (source0.tex lines 474--486). The article's complete original proof of it is delivered with it (source0.tex lines 487--558, one \texttt{proof} environment of 72 lines). No mathematics is written, repaired or re-derived: the statement and the proof are the article's own lines, in the article's own notation, and no retained line is substituted or re-flowed.  Retained uncounted as the source-local context the proof needs: lines 287--296; lines 314--322; lines 325--330.  Nothing was omitted from any retained span, so the package carries no omission record.  No retained line contains a citation, so the wrapper carries no bibliography and no bibliography entry.  No retained line contains a figure environment, an image-inclusion command or drawing code, so no image is delivered and the package declares no asset dependency.  Editorial formatting only, in addition: an AMS \texttt{amsart} preamble that restates the article's own declarations and macros used by the retained lines, namely the environments \texttt{lemma}, \texttt{theorem} on separate counters, together with the standard packages those lines need.  The retained environments print in this document as Lemma 1, Theorem 1 in order of appearance, whereas the article prints the same items with the numbering of its own full document.  The complete native source of the article as distributed in the arXiv e-print for this exact version is bundled unchanged as \texttt{sources/198796/source0.tex}.
 \begin{lemma} \label{lema.1}
 Let $A$ be a Young function. Then
 $$
 \|u*v\|_{L^A(\R^n)} \leq \|v\|_{L^1(\R^n)}\|u\|_{L^A(\R^n)}
 $$
 for every $u\in L^A(\R^n)$ and $v\in L^1(\R^n)$. Moreover,
 $$
 \int_{\R^n}A(|u*v|)\,dx \leq\int_{\R^n} A\left(\|v\|_{L^1(\R^n)} |u|\right)\,dx.
 $$
 \end{lemma}

\begin{lemma} \label{lema.G}
The Bessel kernel satisfies the following properties:
\begin{itemize}
\item[i)] $G_s\in L^1(\mathbb{R}^n)$ for all $s>0$,
\item[ii)] $\nabla G_s\in L^1(\mathbb{R}^n)$ for all $s>1$,
\item[iii)] $G(x) = G(-x)$ for all $x\in \R^n$,
\item[iv)] $G(x)$ is decreasing.
\end{itemize}
\end{lemma}

\begin{lemma} \label{lema.2}
For $s\in(0,1)$, it holds that
$$
\int_{\R^n} |G_s(x+h)-G_s(x)|\,dx \leq \frac{C}{1-s} |h|^s.
$$
\end{lemma}

\begin{theorem} \label{teo.s1}
Let $u\in H^{s',A}(\R^n)$ and $0<s<s'<1$. Then, $u\in W^{s, A}(\R^n)$ and there exists $C>0$ independently of $u$ such that
$$
\|u\|_{W^{s,A}(\R^n)} \leq \frac{C}{s'-s} \|u\|_{H^{s',A}(\R^n)}.
$$
Moreover, if $u=G_{s'}* f$, with $f\in L^A(\R^n)$, then the following modular inequality holds
\begin{align*}
\int_{\R^n} A(|u|)\,dx +&\int_{\R^n}\int_{\R^n}  A\left(\frac{|u(x+h)-u(x)|}{|h|^s}\right)\frac{dxdh}{|h|^n}\leq \\
&\leq  \frac{C_1}{s'-s} \int_{\R^n} A(C_2|f|)\,dx +  \int_{\R^n} A(\|G_{s'}\|_{L^1(\R^n)}|f|)\,dx,
\end{align*}

where $C_1=\frac{2n\omega n}{s}$ and $C_2=2\|G_{s'}\|_{L^1(\R^n)} + \frac{c}{1-s'}$.
\end{theorem}

\begin{proof}
Let $u\in H^{s',A}(\R^n)$ and assume that $u=G_{s'}* f$ for some  $f\in L^A(\R^n)$. Since we can write for $h\neq 0$,
\begin{align*}
\frac{u(x+h)-u(x)}{|h|^s} = \int_{\R^n} (G_{s'}(y+h)-G_{s'}(y))f(x-y)|h|^{-s}\,dy,
\end{align*}
by using Lemma \ref{lema.1} we get that
\begin{align*}
\int_{\R^n} A\left(\frac{|u(x+h)-u(x)|}{|h|^s}\right)\,dx &\leq \int_{\R^n} A\left(\int_{\R^n} |G_{s'}(y+h)-G_{s'}(y))f(x-y)||h|^{-s}\,dy \right)\,dx\\
&\leq \int_{\R^n} A\left( \int_{\R^n}|G_{s'}(y+h)-G_{s'}(y)|\,dy \cdot |f(x)| |h|^{-s}\right)\,dx\\
&= \int_{\R^n} A\left( w_{s'}(h) |f(x)| |h|^{-s}\right)dx,
\end{align*}
where we have denoted $w_{s'}(h):=\int_{\R^n}|G_{s'}(y+h)-G_{s'}(y)|\,dy $. Multiplying by $|h|^{-n}$ and integrating with respect to $h$ gives that
\begin{align} \label{ec.1}
\int_{\R^n}\int_{\R^n} A\left(\frac{|u(x+h)-u(x)|}{|h|^s}\right)\frac{dxdh}{|h|^n}
&\leq \int_{\R^n} \int_{\R^n} A\left( |w_{s'}(h)| |f(x)| |h|^{-s}\right)\frac{dxdh}{|h|^n}.
\end{align}
We will estimate the following integral:
\begin{align*}
\int_{\R^n} A\left( |w_{s'}(h)| |f(x)| |h|^{-s}\right)|h|^{-n}dh = \left(\int_{|h|\leq 1} + \int_{|h|>1} \right) A\left( |w_{s'}(h)| |f(x)| |h|^{-s}\right)|h|^{-n}dh:=I_1+ I_2.
\end{align*}

\underline{Case $|h|> 1$}. Using Lemma \ref{lema.G} we have that $|w_{s'}(h)|\leq 2\|G_{s'}\|_{L^1(\R^n)}$, so,
\begin{align*}
I_2&= \int_{|h|>1} A\left( |w_{s'}(h)| |f(x)| |h|^{-s}\right)|h|^{-n}dh\\
&\leq 
\int_{|h|>1} A\left( 2\|G_{s'}\|_{L^1(\R^n)}|f(x)| |h|^{-s}\right)|h|^{-n}dh \\
&=
\int_1^\infty  A\left( 2\|G_{s'}\|_{L^1(\R^n)} |f(x)| r^{-s}\right)r^{-1}dh\\
&\leq 
A\left(2\|G_{s'}\|_{L^1(\R^n)} |f(x)| \right) n\omega_n\int_1^\infty  r^{-1-s}dh=  \frac{n\omega_n}{s} A\left(2\|G_{s'}\|_{L^1(\R^n)} |f(x)| \right),
\end{align*}
since $r^{-s}<1$ and $A$ is convex. 

\underline{Case $|h|\leq 1$}. By Lemma \ref{lema.2}, $|w_{s'}(h)|\leq \frac{C}{1-s'}|h|^{s'}$, so, using the convexity of $A$ gives that
\begin{align*}
I_1&\leq \int_{|h|\leq 1} A\left( \tfrac{C}{1-s}|h|^{s'} |f(x)| |h|^{-s}\right)|h|^{-n}dh\\
&\leq A(\tfrac{C}{1-s'}|f(x)|) \int_{0}^1 |h|^{-n+s'-s}dh\leq \frac{n\omega_n}{s'-s}A(\tfrac{C}{1-s'}|f(x)|) 
\end{align*}
Replacing the bound for $I_1$ and $I_2$ into \eqref{ec.1} gives that
$$
\int_{\R^n}\int_{\R^n} A\left(\frac{|u(x+h)-u(x)|}{|h|^s}\right)\frac{dxdh}{|h|^n}
\leq  \frac{C_1}{s'-s} \int_{\R^n} A(C_2 |f(x)|)\,dx
$$
where $C_1=\frac{2n\omega n}{s}$ and $C_2=2\|G_{s'}\|_{L^1(\R^n)} + \frac{c}{1-s'}$.

  
Finally, let $\lambda := C C_2\|f\|_{L^A(\R^n)}$ for some $C>1$ such that $C>\frac{C_1}{s'-s}$. Observe that a similar computation gives that
$$
\int_{\R^n}\int_{\R^n} A\left(\frac{|u(x+h)-u(x)|}{\lambda |h|^s}\right)\frac{dxdh}{|h|^n}
\leq  \frac{C_1}{s'-s} \int_{\R^n} A\left(\frac{C_2|f(x)|}{\lambda}\right)\,dx,
$$
but, since $C>1$
$$
\frac{C_1}{s'-s} \int_{\R^n} A\left(\frac{C_2|f(x)|}{\lambda}\right)\,dx 
\leq 
\frac{C_1}{s'-s} \frac{1}{C}\int_{\R^n} A\left(\frac{|f(x)|}{ \|f\|_{L^A(\R^n)}}\right)\,dx\leq 1
$$
by definition of the Luxemburg norm. Then, using again the definition of the norm
$$
[u]_{W^{s,A}(\R^n)} \leq CC_2 \|f\|_{L^A(\R^n)}= CC_2 \|u\|_{H^{s,A}(\R^n)}.
$$

Finally, observe that by Lemma \ref{lema.1}
$$
\|u\|_{L^A(\R^n)}=\|G_s* f\|_{L^A(\R^n)}\leq  \|G_s\|_{L^1(\R^n)}\|f\|_{L^A(\R^n)} = \|G_s\|_{L^1(\R^n)}\|u\|_{H^{s,A}(\R^n)}
$$
and
$$
\int_{\R^n} A(|G_s*f|)\,dx \leq \int_{\R^n} A(\|G_s\|_{L^1(\R^n)}|f|)\,dx.
$$
This concludes the proof.
\end{proof}

\end{document}
